\documentclass[a4paper,12pt]{article}
\usepackage[utf8]{inputenc}
\usepackage[T1]{fontenc}
\usepackage[ngerman]{babel}
\usepackage{amsmath,amssymb}
\usepackage{booktabs}
\usepackage{enumitem}
\usepackage[margin=2.5cm]{geometry}
\usepackage{tcolorbox}

\tcbuselibrary{breakable}

\newtcolorbox{merkbox}[1][]{
  colback=blue!5,
  colframe=blue!40,
  fonttitle=\bfseries,
  title={Merke},
  #1
}

\newtcolorbox{analogie}[1][]{
  colback=green!5,
  colframe=green!40,
  fonttitle=\bfseries,
  title={Analogie},
  #1
}

\begin{document}

\begin{center}
{\LARGE\bfseries \"Ubungsblatt 2 -- Erkl\"arungen}\\[0.5em]
{\large VC Logik}
\end{center}

\bigskip

%% ============================================================
%% AUFGABE 1
%% ============================================================
\section{Aufgabe 1: Kaffeewelt}

\subsection{Was wird hier verlangt?}

Wir haben ein Geb\"aude mit R\"aumen. In manchen R\"aumen riecht
es nach Kaffee (Variable~$G$), in manchen steht eine Kaffeetasse
(Variable~$K$). Die Wissensbasis~KB fasst unsere Beobachtungen
und Regeln zusammen. Wir sollen \emph{alle} m\"oglichen
Belegungen der Variablen (Interpretationen) auflisten und
herausfinden, welche davon mit KB vertr\"aglich sind (Modelle).

\begin{analogie}
Stell dir vor, du bist ein Detektiv. Du betrittst Raum~$(1{,}1)$
und riechst \emph{keinen} Kaffee. Dann gehst du in Raum~$(1{,}2)$
und riechst Kaffee. Jetzt willst du herausfinden, in welchen
der angrenzenden R\"aume $(2{,}1)$, $(2{,}2)$, $(1{,}3)$ eine
Kaffeetasse stehen k\"onnte -- das sind die \emph{Modelle}.
\end{analogie}

\subsection{Was kodiert die Wissensbasis?}

\begin{align*}
\text{KB} &\equiv \underbrace{\neg G_{1,1}}_{\text{kein Geruch in }(1,1)}
  \;\wedge\; \underbrace{G_{1,2}}_{\text{Geruch in }(1,2)}
  \;\wedge \\
& \underbrace{(K_{2,1} \to G_{1,1})}_{\text{Tasse in }(2,1)
  \Rightarrow\text{ Geruch in }(1,1)}
  \;\wedge\;
  \underbrace{(K_{2,2} \to G_{1,2})}_{\text{Tasse in }(2,2)
  \Rightarrow\text{ Geruch in }(1,2)}
  \;\wedge\;
  \underbrace{(K_{1,3} \to G_{1,2})}_{\text{Tasse in }(1,3)
  \Rightarrow\text{ Geruch in }(1,2)}
  \;\wedge \\
& \underbrace{(G_{1,1} \to K_{2,1})}_{\text{Geruch in }(1,1)
  \Rightarrow\text{ Tasse in }(2,1)}
  \;\wedge\;
  \underbrace{(G_{1,2} \to K_{2,2} \vee K_{1,3})}_{\text{Geruch in }(1,2)
  \Rightarrow\text{ mind.\ eine Tasse}}
\end{align*}

Zusammen ergeben die letzten f\"unf Konjunkte zwei
Bikonditionale ("`genau dann, wenn"'):
\begin{itemize}
\item $G_{1,1} \leftrightarrow K_{2,1}$
      \quad (Geruch in $(1{,}1)$ $\Leftrightarrow$ Tasse in $(2{,}1)$)
\item $G_{1,2} \leftrightarrow K_{2,2} \vee K_{1,3}$
      \quad (Geruch in $(1{,}2)$ $\Leftrightarrow$ Tasse in $(2{,}2)$ oder $(1{,}3)$)
\end{itemize}

\subsection{Schritt f\"ur Schritt: Constraints ableiten}

\begin{enumerate}[leftmargin=2cm]
\item[\textbf{Schritt 1:}]
\textbf{Beobachtungen auswerten.}

$\neg G_{1,1}$ erzwingt $G_{1,1} = 0$.
$G_{1,2}$ erzwingt $G_{1,2} = 1$.

\item[\textbf{Schritt 2:}]
\textbf{Implikationen mit festen Werten auswerten.}

$K_{2,1} \to G_{1,1}$: Da $G_{1,1} = 0$, muss $K_{2,1} = 0$
sein (sonst w\"are $1 \to 0 = 0$).

$K_{2,2} \to G_{1,2}$: Da $G_{1,2} = 1$, ist dies immer wahr
($K_{2,2} \to 1 = 1$). Keine Einschr\"ankung f\"ur $K_{2,2}$.

$K_{1,3} \to G_{1,2}$: Ebenso immer wahr.

$G_{1,1} \to K_{2,1}$: $0 \to K_{2,1} = 1$, immer wahr.

$G_{1,2} \to K_{2,2} \vee K_{1,3}$: Da $G_{1,2} = 1$,
brauchen wir $K_{2,2} \vee K_{1,3} = 1$, also mindestens
eine Tasse in $(2{,}2)$ oder $(1{,}3)$.

\item[\textbf{Schritt 3:}]
\textbf{Zusammenfassung.}

KB ist genau dann wahr, wenn:
$G_{1,1}=0$, \;$G_{1,2}=1$, \;$K_{2,1}=0$, \;und
$K_{2,2} \vee K_{1,3} = 1$.
\end{enumerate}

\begin{merkbox}
Von den $2^5 = 32$ Interpretationen sind genau \textbf{3~Modelle}:

\smallskip
\begin{tabular}{cccccc}
$G_{1,1}$ & $G_{1,2}$ & $K_{2,1}$ & $K_{2,2}$ & $K_{1,3}$ & Bedeutung \\
\midrule
0 & 1 & 0 & 0 & 1 & Tasse nur in $(1{,}3)$ \\
0 & 1 & 0 & 1 & 0 & Tasse nur in $(2{,}2)$ \\
0 & 1 & 0 & 1 & 1 & Tassen in $(2{,}2)$ und $(1{,}3)$ \\
\end{tabular}
\end{merkbox}

\subsection{Die 8~Welten aus der Vorlesung}

Die 8~m\"oglichen Welten entsprechen den $2^3 = 8$ Kombinationen
der Tassenvariablen $K_{2,1}$, $K_{2,2}$, $K_{1,3}$.
In jeder Welt sind die Geruchswerte durch die Bikonditionale
eindeutig bestimmt:

\begin{center}
\renewcommand{\arraystretch}{1.25}
\begin{tabular}{c|ccc|cc|c|l}
\toprule
Welt & $K_{2,1}$ & $K_{2,2}$ & $K_{1,3}$
     & $G_{1,1}$ & $G_{1,2}$ & Modell? & Grund \\
\midrule
$W_1$ & 0 & 0 & 0 & 0 & 0 & Nein & $G_{1,2} \neq 1$ \\
$W_2$ & 0 & 0 & 1 & 0 & 1 & \textbf{Ja} & \\
$W_3$ & 0 & 1 & 0 & 0 & 1 & \textbf{Ja} & \\
$W_4$ & 0 & 1 & 1 & 0 & 1 & \textbf{Ja} & \\
$W_5$ & 1 & 0 & 0 & 1 & 0 & Nein & $G_{1,1} \neq 0$ \\
$W_6$ & 1 & 0 & 1 & 1 & 1 & Nein & $G_{1,1} \neq 0$ \\
$W_7$ & 1 & 1 & 0 & 1 & 1 & Nein & $G_{1,1} \neq 0$ \\
$W_8$ & 1 & 1 & 1 & 1 & 1 & Nein & $G_{1,1} \neq 0$ \\
\bottomrule
\end{tabular}
\end{center}

Die restlichen $32 - 8 = 24$ Interpretationen haben
Geruchswerte, die \emph{nicht} zu den Tassen passen
(z.\,B.\ Geruch ohne Tasse oder Tasse ohne Geruch) --
sie verletzen die Bikonditionale in~KB und sind daher
ebenfalls keine Modelle.

\newpage

%% ============================================================
%% AUFGABE 2
%% ============================================================
\section{Aufgabe 2: G\"ultigkeit, Erf\"ullbarkeit, \"Aquivalenz,
logische Konsequenz}

\subsection{Die wichtigsten Begriffe}

\begin{merkbox}
\begin{tabular}{lp{10cm}}
\textbf{G\"ultig} (Tautologie): &
Formel ist in \emph{allen} Interpretationen wahr. \\[3pt]
\textbf{Erf\"ullbar}: &
Formel ist in \emph{mindestens einer} Interpretation wahr. \\[3pt]
\textbf{\"Aquivalent} ($F_1 \equiv F_2$): &
Beide Formeln haben in \emph{jeder} Interpretation
denselben Wahrheitswert. \\[3pt]
\textbf{Log.\ Konsequenz} ($F_1 \models F_2$): &
\emph{Jedes} Modell von $F_1$ ist auch Modell von $F_2$.
\end{tabular}
\end{merkbox}

\begin{analogie}
\begin{itemize}
\item \textbf{G\"ultig} = ein Gesetz, das \emph{immer} gilt,
      egal was passiert.
\item \textbf{Erf\"ullbar} = es gibt \emph{mindestens eine}
      Situation, in der es stimmt.
\item \textbf{\"Aquivalent} = zwei Aussagen, die immer dasselbe
      sagen (wie "`Es regnet"' und "`Wasser f\"allt vom Himmel"').
\item \textbf{Konsequenz} = wenn das eine gilt, \emph{muss}
      auch das andere gelten.
\end{itemize}
\end{analogie}

\subsection{Wahrheitstabelle}

\begin{center}
\renewcommand{\arraystretch}{1.25}
\begin{tabular}{cc|ccccc}
\toprule
$P$ & $Q$
  & (a)\;$P$
  & (b)\;$Q$
  & (c)\;$P \wedge Q$
  & (d)\;$\neg P \vee \neg Q$
  & (e)\;$\neg(\neg P \vee \neg Q)$ \\
\midrule
0 & 0 & 0 & 0 & 0 & 1 & 0 \\
0 & 1 & 0 & 1 & 0 & 1 & 0 \\
1 & 0 & 1 & 0 & 0 & 1 & 0 \\
1 & 1 & 1 & 1 & 1 & 0 & 1 \\
\bottomrule
\end{tabular}
\end{center}

\subsection{Analyse}

\subsubsection*{G\"ultigkeit}

Eine Formel ist g\"ultig, wenn ihre Spalte \emph{nur} Einsen
enth\"alt. Das trifft auf \textbf{keine} der f\"unf Formeln zu.

\subsubsection*{Erf\"ullbarkeit}

Eine Formel ist erf\"ullbar, wenn ihre Spalte \emph{mindestens
eine} Eins enth\"alt. Das trifft auf \textbf{alle f\"unf} Formeln
zu.

\subsubsection*{\"Aquivalenz}

Zwei Formeln sind \"aquivalent, wenn ihre Spalten
\emph{identisch} sind.

\begin{itemize}
\item (c) hat die Werte $(0,0,0,1)$ und (e) hat die Werte
$(0,0,0,1)$ $\Rightarrow$ \textbf{(c) $\equiv$ (e)}.

Dies ist das \textbf{Gesetz von De~Morgan}:
$P \wedge Q \equiv \neg(\neg P \vee \neg Q)$.

\item Alle anderen Paare haben unterschiedliche Spalten
$\Rightarrow$ nicht \"aquivalent.
\end{itemize}

\begin{merkbox}
\textbf{Zus\"atzliche Beobachtung:} (d) hat die Werte
$(1,1,1,0)$ -- das ist das genaue Gegenteil von (c) bzw.\ (e).
Also gilt: $(d) \equiv \neg(c) \equiv \neg(e)$. Das ist ebenfalls
De~Morgan: $\neg P \vee \neg Q \equiv \neg(P \wedge Q)$.
\end{merkbox}

\subsubsection*{Logische Konsequenz}

$F_1 \models F_2$ gilt, wenn $\text{Mod}(F_1) \subseteq \text{Mod}(F_2)$.

\smallskip
Die Modellmengen (als Mengen von $(P,Q)$-Paaren):
\begin{align*}
\text{Mod}(a) &= \{(1,0),\;(1,1)\} \\
\text{Mod}(b) &= \{(0,1),\;(1,1)\} \\
\text{Mod}(c) &= \{(1,1)\} \\
\text{Mod}(d) &= \{(0,0),\;(0,1),\;(1,0)\} \\
\text{Mod}(e) &= \{(1,1)\}
\end{align*}

\textbf{Konsequenzen, die gelten:}
\begin{itemize}
\item $(c) \models (a)$: \;$\{(1,1)\} \subseteq \{(1,0),(1,1)\}$
      \;\checkmark\quad
      ($P \wedge Q$ impliziert $P$)
\item $(c) \models (b)$: \;$\{(1,1)\} \subseteq \{(0,1),(1,1)\}$
      \;\checkmark\quad
      ($P \wedge Q$ impliziert $Q$)
\item $(e) \models (a)$ und $(e) \models (b)$: \;\checkmark\quad
      (folgt aus $(c) \equiv (e)$)
\end{itemize}

\begin{analogie}
\textbf{Warum $(c) \models (a)$?}
Wenn du wei\ss t, dass \emph{sowohl} $P$ als auch $Q$ wahr sind
($P \wedge Q$), dann ist insbesondere $P$ wahr.
Das Spezielle impliziert das Allgemeine.
\end{analogie}

\textbf{Beispiel, warum $(a) \not\models (c)$:}
$(1,0)$ ist ein Modell von $P$, aber kein Modell von $P \wedge Q$
(da $Q=0$). Also impliziert $P$ allein \emph{nicht} $P \wedge Q$.

\newpage

%% ============================================================
%% AUFGABE 3
%% ============================================================
\section{Aufgabe 3: Erf\"ullbarkeitstester}

\subsection{Was ist ein Erf\"ullbarkeitstester?}

\begin{analogie}
Ein Erf\"ullbarkeitstester (SAT-Solver) ist wie ein
\textbf{Automat}: Du gibst ihm eine Formel, und er antwortet
mit "`erf\"ullbar"' oder "`unerf\"ullbar"'. Mehr kann er nicht --
aber das reicht, um viele Fragen zu beantworten!
\end{analogie}

\subsection{Die Schl\"usselbeziehungen}

\begin{merkbox}
\renewcommand{\arraystretch}{1.4}
\begin{tabular}{lll}
\textbf{Frage} & \textbf{Gib dem Tester} & \textbf{Antwort} \\
\midrule
Ist $F$ g\"ultig? & $\neg F$ & unerf.\ $\Rightarrow$ ja \\
Ist $F$ Kontradiktion? & $F$ & unerf.\ $\Rightarrow$ ja \\
Folgt $F_2$ aus $F_1$? & $F_1 \wedge \neg F_2$
  & unerf.\ $\Rightarrow$ ja \\
Ist $F$ Tautologie? & $\neg F$ & unerf.\ $\Rightarrow$ ja \\
\end{tabular}
\end{merkbox}

\textbf{Idee dahinter:} Wenn die Negation einer Aussage
unerf\"ullbar ist (es gibt kein Gegenbeispiel), dann muss
die Aussage selbst immer wahr sein.

\subsection{Umformulierungen}

\begin{enumerate}[(a)]
\item \textbf{Ist $(P \vee (\neg P \to Q)) \leftrightarrow (P \vee Q)$
      g\"ultig?}

      \textbf{Umformulierung:} Ist
      $\neg\bigl((P \vee (\neg P \to Q)) \leftrightarrow (P \vee Q)\bigr)$
      unerf\"ullbar?

      Falls ja $\Rightarrow$ die Formel ist g\"ultig.

      \smallskip
      \textit{Antwort (Flei\ss aufgabe):}
      $\neg P \to Q \equiv P \vee Q$, also vereinfacht sich
      die Formel zu $(P \vee Q) \leftrightarrow (P \vee Q)$,
      was stets wahr ist. Die Negation ist unerf\"ullbar
      $\Rightarrow$ \textbf{g\"ultig}.

\item \textbf{Folgt $P \to Q$ aus $\neg Q \to \neg P$?}

      \textbf{Umformulierung:} Ist
      $(\neg Q \to \neg P) \wedge \neg(P \to Q)$
      unerf\"ullbar?

      Falls ja $\Rightarrow$ es folgt.

      \smallskip
      \textit{Antwort (Flei\ss aufgabe):}
      $\neg Q \to \neg P$ ist die Kontraposition von $P \to Q$.
      Beide sind \"aquivalent, also folgt $P \to Q$ aus
      $\neg Q \to \neg P$ $\Rightarrow$ \textbf{ja, es folgt}.

\item \textbf{Ist $P \leftrightarrow Q \wedge P$ eine Kontradiktion?}

      Vollst\"andig geklammert: $P \leftrightarrow (Q \wedge P)$.

      \textbf{Umformulierung:} Ist
      $P \leftrightarrow (Q \wedge P)$ unerf\"ullbar?

      Falls ja $\Rightarrow$ Kontradiktion.

      \smallskip
      \textit{Antwort (Flei\ss aufgabe):}
      F\"ur $P=0$, $Q=0$: $0 \leftrightarrow (0 \wedge 0) =
      0 \leftrightarrow 0 = 1$.
      Die Formel ist erf\"ullbar $\Rightarrow$
      \textbf{keine Kontradiktion}.

\item \textbf{Ist $P \leftrightarrow P \vee \bot$ eine Tautologie?}

      Vollst\"andig geklammert: $P \leftrightarrow (P \vee \bot)$.

      \textbf{Umformulierung:} Ist
      $\neg\bigl(P \leftrightarrow (P \vee \bot)\bigr)$ unerf\"ullbar?

      Falls ja $\Rightarrow$ Tautologie.

      \smallskip
      \textit{Antwort (Flei\ss aufgabe):}
      $P \vee \bot = P \vee 0 = P$, also wird die Formel zu
      $P \leftrightarrow P$, was stets wahr ist.
      Die Negation ist unerf\"ullbar $\Rightarrow$
      \textbf{Tautologie}.
\end{enumerate}

\end{document}
